\chapter{Motivación de la Elección de Software}

La elección del software empleado responde a criterios de fiabilidad, adecuación al modelo relacional y eficiencia en el desarrollo. Como Sistema Gestor de Base de Datos (SGBD) se ha optado por un entorno relacional basado en Oracle, accedido mediante \textit{SQL Developer}. Esta elección se justifica por su madurez tecnológica y por ofrecer un soporte sólido de transacciones conforme a propiedades \textit{ACID}, garantizando consistencia e integridad de los datos en operaciones concurrentes. Asimismo, permite trabajar de forma natural con un diseño relacional normalizado, aplicando restricciones (\textit{PRIMARY KEY}, \textit{FOREIGN KEY}, \textit{CHECK}, \textit{NOT NULL}) y asegurando la integridad referencial, aspecto esencial para preservar la coherencia del modelo de datos. Adicionalmente, el motor incorpora mecanismos de optimización de consultas y herramientas avanzadas para depuración y administración, lo que resulta especialmente útil en la verificación del rendimiento y la trazabilidad de errores.

Por su parte, \textit{SQL Developer} se ha utilizado como herramienta principal de modelado, ejecución y depuración de sentencias SQL por su interfaz gráfica, su estabilidad, su amplia documentación y su uso extendido en contextos académicos y profesionales. Finalmente, se ha empleado \textit{Python} como soporte complementario para automatizar procesos, realizar tareas de preprocesamiento y facilitar el análisis de datos, aprovechando su ecosistema de bibliotecas y su integración con bases de datos mediante conectores estándar. Esta combinación proporciona un flujo de trabajo reproducible, eficiente y alineado con prácticas habituales en proyectos de ingeniería de datos.
